\documentclass[twocolumn]{article}
\usepackage{booktabs}
\usepackage{multirow}
\usepackage{threeparttable}
\usepackage{graphicx}
\usepackage{caption}
\usepackage{subcaption}
\usepackage{amsmath}

\title{Binary Spine X-Ray Classification:\\
  Classical CNN vs.\ Discrete-Variable and Continuous-Variable Quantum Neural Networks}
\author{QStrata Research Team}
\date{2026}

\begin{document}
\maketitle

\section{Experiment Setup}
% Experiment Setup Table
% Generated by Q60-PAPER-ARTIFACTS-EXPORT

\begin{table}[htbp]
\centering
\caption{Unified experiment setup for all binary classification benchmarks (Q56–Q58).
  All three experiments share the same training protocol; only the classification head differs.}
\label{tab:experiment_setup}
\begin{tabular}{lp{8.5cm}}
\toprule
\textbf{Setting} & \textbf{Value} \\
\midrule
\multicolumn{2}{l}{\textit{Datasets}} \\
VinDr-SpineXR          & Real clinical DICOM X-rays; binary: abnormal vs.\ normal; \\
                        & train / val / test split per official partitions \\
BUU-LSPINE             & Synthetic lumbar spine X-rays; binary classification \\
\midrule
\multicolumn{2}{l}{\textit{Backbone}} \\
Architecture            & MobileNetV3-Large (ImageNet-pretrained, all layers frozen) \\
Feature dimension       & 960-dimensional pooled feature vector \\
\midrule
\multicolumn{2}{l}{\textit{Training protocol}} \\
Epochs                  & 4 \\
Batch size              & 32 \\
Optimiser               & Adam \\
Learning rate           & $1 \times 10^{-3}$ \\
Loss                    & Binary cross-entropy with logits \\
Seeds                   & 42, 7, 123 (3 independent runs per experiment) \\
\midrule
\multicolumn{2}{l}{\textit{Evaluation}} \\
Primary metric          & AUROC (area under ROC curve) \\
Secondary metrics       & F1 score, accuracy \\
Confidence intervals    & 95\% CI computed over 3 seeds \\
Hardware                & Single CPU core (no GPU) \\
Framework               & PyTorch + custom \texttt{qcore} DV/CV simulators (no external quantum libraries) \\
\midrule
\multicolumn{2}{l}{\textit{Quantum circuit settings (QNN only)}} \\
DV-QNN qubits           & 4 qubits, depth 1, $\alpha = 0.1$, \texttt{medical\_ansatz} \\
CV-QNN modes            & 2 modes, depth 1, squeezing cap $= 1.5$, \texttt{GVA} ansatz \\
\bottomrule
\end{tabular}
\end{table}


\section{Architecture Comparison}
% Architecture Comparison Table
% Generated by Q60-PAPER-ARTIFACTS-EXPORT

\begin{table}[htbp]
\centering
\caption{Architecture comparison across the three binary classification pipelines.
  All three share a frozen MobileNetV3-Large backbone; only the head/circuit differs.}
\label{tab:architecture}
\begin{tabular}{llllll}
\toprule
\textbf{Model} & \textbf{Backbone} & \textbf{Head / Circuit} & \textbf{Simulator} & \textbf{Layers / Depth} & \textbf{Activation} \\
\midrule
Classical CNN (Q56)
  & MobileNetV3-Large (frozen) + Linear(960$\to$128, frozen proj.)
  & 3-layer MLP
  & ---
  & FC(128$\to$16), FC(16$\to$8), FC(8$\to$2)
  & ReLU + Softmax \\
DV-QNN (Q57)
  & MobileNetV3-Large (frozen)
  & 4-qubit VQC
  & qcore DV (discrete)
  & depth$=1$, $\alpha{=}0.1$
  & Pauli-Z measurement \\
CV-QNN (Q58)
  & MobileNetV3-Large (frozen)
  & 2-mode CV circuit
  & qcore Gaussian (CV)
  & depth$=1$, squeeze cap$=1.5$
  & Homodyne measurement \\
\bottomrule
\end{tabular}
\begin{tablenotes}
\small
\item All backbones initialised from ImageNet-pretrained weights; backbone parameters are frozen during training.
\item DV-QNN uses a custom \texttt{medical\_ansatz}; CV-QNN uses \texttt{GaussianBackend} + \texttt{GVA} ansatz.
\item MLP head is the Q34A reference head used across all Phase 6b benchmarks.
\end{tablenotes}
\end{table}


\section{Parameter Efficiency}
% Parameter Count and Efficiency Table
% Generated by Q60-PAPER-ARTIFACTS-EXPORT

\begin{table}[htbp]
\centering
\caption{Trainable parameter counts and efficiency metrics.
  The backbone (MobileNetV3-Large) contributes 5,606,040 frozen parameters in all cases.
  Efficiency ratio is the AUROC achieved on VinDr-SpineXR relative to Classical CNN per trainable parameter.}
\label{tab:parameters}
\begin{tabular}{lrrrrr}
\toprule
\textbf{Model}
  & \textbf{Trainable}
  & \textbf{Total}
  & \textbf{Frozen}
  & \textbf{Train. \%}
  & \textbf{AUROC / 1k params} \\
\midrule
Classical CNN (Q56) & 2{,}250  & 5{,}608{,}290 & 5{,}606{,}040 & 0.040\% & 432.9 \\
DV-QNN (Q57)        & 574      & 5{,}606{,}614 & 5{,}606{,}040 & 0.010\% & 154.0 \\
CV-QNN (Q58)        & 532      & 5{,}606{,}572 & 5{,}606{,}040 & 0.009\% & 179.2 \\
\bottomrule
\end{tabular}
\begin{tablenotes}
\small
\item CV-QNN achieves 98.0\% of Classical CNN's VinDr AUROC with 76.4\% fewer trainable parameters (532 vs.\ 2{,}250).
\item DV-QNN achieves 90.9\% of Classical CNN's VinDr AUROC with 74.5\% fewer trainable parameters (574 vs.\ 2{,}250).
\item AUROC / 1k params $= 1000 \times \text{AUROC} / \text{trainable params}$; higher is more parameter-efficient.
\end{tablenotes}
\end{table}


\section{Results}
% Main Results Table — Binary Classification Benchmark
% Q56 (Classical CNN), Q57 (DV-QNN), Q58 (CV-QNN) on VinDr-SpineXR and BUU-LSPINE
% Generated by Q60-PAPER-ARTIFACTS-EXPORT

\begin{table}[htbp]
\centering
\caption{Binary classification performance across architectures and datasets (mean $\pm$ std over 3 seeds).
  Best result per column in \textbf{bold}. $\Delta$ columns show difference relative to Classical CNN baseline.
  Runtime measured on a single CPU core, averaged over 3 seeds.}
\label{tab:main_results}
\resizebox{\textwidth}{!}{%
\begin{tabular}{llcccccc}
\toprule
\multirow{2}{*}{\textbf{Architecture}} &
\multirow{2}{*}{\textbf{Method}} &
\multicolumn{3}{c}{\textbf{VinDr-SpineXR (Real DICOM)}} &
\multicolumn{2}{c}{\textbf{BUU-LSPINE (Synthetic)}} &
\multirow{2}{*}{\textbf{Runtime (s)}} \\
\cmidrule(lr){3-5}\cmidrule(lr){6-7}
& & \textbf{AUROC} & \textbf{F1} & $\Delta$\textbf{AUROC} & \textbf{AUROC} & \textbf{F1} & \\
\midrule
MobileNetV3-L + MLP         & Classical CNN  & \textbf{0.9731 $\pm$ 0.0018} & \textbf{0.9153} & ---              & \textbf{1.0000}          & \textbf{1.0000} & \textbf{0.57}  \\
MobileNetV3-L + DV-QNN      & DV-QNN (Q57)   & 0.8842 $\pm$ 0.0008          & 0.7888          & $-$0.0889        & \textbf{1.0000}          & 0.9898          & 263.50         \\
MobileNetV3-L + CV-QNN      & CV-QNN (Q58)   & 0.9534 $\pm$ 0.0010          & 0.8791          & $-$0.0197        & \textbf{1.0000}          & \textbf{1.0000} & 56.20          \\
\bottomrule
\end{tabular}%
}
\begin{tablenotes}
\small
\item AUROC CI$_{95}$: CNN [0.9685, 0.9776]; DV-QNN [0.8822, 0.8863]; CV-QNN [0.9510, 0.9558].
\item Runtime is total training + inference time per run on VinDr-SpineXR.
\item All models trained for 4 epochs with LR $= 10^{-3}$; seeds $\{42, 7, 123\}$.
\end{tablenotes}
\end{table}


\section{Figures}

\begin{figure}[htbp]
  \centering
  \includegraphics[width=\linewidth]{figures/dataset_overview}
  \caption{Dataset overview: VinDr-SpineXR (real DICOM) and BUU-LSPINE (synthetic).}
  \label{fig:datasets}
\end{figure}

\begin{figure*}[htbp]
  \centering
  \includegraphics[width=\linewidth]{figures/main_comparison_figure}
  \caption{Five-panel comparison across architectures and datasets.
    (A) VinDr AUROC. (B) BUU-LSPINE AUROC.
    (C) Training runtime (log scale). (D) Cross-dataset generalisation gap.
    (E) Trainable parameter counts.}
  \label{fig:main_comparison}
\end{figure*}

\end{document}
